% ============================================================================
% BOLUM 6: SONUC VE DEGERLENDIRME
% ============================================================================
% Kurumsal Temizlik Takip ve Yonetim Sistemi (KTTYS)
% Muhammet Anil YAGIZ - 213255046
% Kirikkale Universitesi - Bilgisayar Muhendisligi Bolumu
% ============================================================================

\chapter{SONUC VE DEGERLENDIRME}
\label{chap:sonuc}

\section{Projenin Ozeti}
\label{sec:proje-ozeti}

Bu bitirme projesi kapsaminda, Kirikkale Universitesinin temizlik 
hizmetlerinin dijitallestirilmesi ve optimize edilmesi amaciyla 
\textbf{Kurumsal Temizlik Takip ve Yonetim Sistemi (KTTYS)} 
gelistirilmistir.

Gelistirilen sistem, modern web teknolojileri kullanilarak 
Progressive Web Application (PWA) mimarisi uzerine insa edilmistir. 
Backend tarafinda Node.js ve Express.js, veritabani olarak SQLite, 
frontend tarafinda HTML5, CSS3 ve JavaScript kullanilmistir.

\section{Elde Edilen Sonuclar}
\label{sec:elde-edilen-sonuclar}

Proje kapsaminda gerceklestirilen calismalar sonucunda asagidaki 
somut ciktilar elde edilmistir:

\subsection{Teknik Ciktilar}
\label{subsec:teknik-ciktilar}

\begin{enumerate}
    \item \textbf{Tam Fonksiyonel Web Uygulamasi:} RESTful API 
    mimarisi uzerine kurulu, 7 adet route modulu ve 30'dan fazla 
    API endpoint iceren calisir durumda bir web uygulamasi.
    
    \item \textbf{Veritabani Tasarimi:} Iliskisel veritabani 
    prensiplerine uygun, 7 adet tablo ve gerekli indekslerden 
    olusan optimize edilmis bir veritabani semasi.
    
    \item \textbf{Guvenlik Altyapisi:} bcrypt ile sifre hashleme, 
    session tabanli kimlik dogrulama ve RBAC ile yetkilendirme.
    
    \item \textbf{Responsive Arayuz:} Masaustu ve mobil cihazlarda 
    uyumlu calisan, kullanici dostu bir web arayuzu.
    
    \item \textbf{PWA Destegi:} Service Worker ve manifest dosyasi 
    ile PWA standartlarina uygun yapilandirma.
\end{enumerate}

\subsection{Islevsel Ciktilar}
\label{subsec:islevsel-ciktilar}

Sistem, asagidaki islevsel yetenekleri basariyla gerceklestirmektedir:

\begin{enumerate}
    \item \textbf{Kullanici Yonetimi:} Farkli rollerde kullanici 
    olusturma, duzenleme ve yetkilendirme.
    
    \item \textbf{Organizasyon Yapisi:} Fakulte, Bina, Departman, 
    Lokasyon hiyerarsisi tanimlama.
    
    \item \textbf{Program Planlama:} Haftalik temizlik programlari 
    olusturma ve personele atama.
    
    \item \textbf{Gorev Yonetimi:} Otomatik gorev olusturma, baslatma, 
    tamamlama, onaylama ve reddetme islemleri.
    
    \item \textbf{Tatil Yonetimi:} Turkiye resmi tatillerinin 
    tanimlanmasi ve gorev sistemine entegrasyonu.
\end{enumerate}

\section{Basari Kriterleri Degerlendirmesi}
\label{sec:basari-degerlendirme}

\begin{table}[H]
\centering
\caption{Basari Kriterleri Degerlendirme Sonuclari}
\label{tab:basari-sonuclari}
\begin{tabular}{|p{3.5cm}|p{2.5cm}|p{2.5cm}|p{2cm}|}
\hline
\textbf{Kriter} & \textbf{Hedef} & \textbf{Sonuc} & \textbf{Durum} \\
\hline
Gorev tamamlama orani & 95\%+ & Test verileri & Basarili \\
\hline
Raporlama suresi & $<$ 1 dakika & $<$ 500ms & Basarili \\
\hline
Sayfa yuklenme suresi & $<$ 3 saniye & $<$ 2 saniye & Basarili \\
\hline
API yanit suresi & $<$ 500ms & $<$ 200ms & Basarili \\
\hline
\end{tabular}
\end{table}

\section{Karsilasilan Zorluklar}
\label{sec:zorluklar}

\subsection{Teknik Zorluklar}
\label{subsec:teknik-zorluklar}

\begin{enumerate}
    \item \textbf{Native Modul Derleme:} Ilk basta kullanilan 
    better-sqlite3 kutuphanesi, Node.js versiyonu ile uyumsuzluk 
    nedeniyle native modul derleme hatasi vermistir. Bu sorun, 
    sql.js kutuphanesine gecilerek cozulmustur.
    
    \item \textbf{Transaction Destegi:} sql.js kutuphanesinin 
    better-sqlite3 ile tam API uyumlulugu saglamak icin ozel 
    bir wrapper katmani gelistirilmistir.
    
    \item \textbf{Asenkron/Senkron Uyum:} sql.js'in asenkron 
    baslatma yapisi ile Express route'larinin senkron beklentisi 
    arasinda uyum saglanmistir.
\end{enumerate}

\section{Ogrenilen Dersler}
\label{sec:ogrenilen-dersler}

Bu proje surecinde edinilen onemli deneyimler sunlardir:

\begin{enumerate}
    \item \textbf{Teknoloji Seciminin Onemi:} Native moduller yerine 
    pure JavaScript cozumlerin secilmesi, platform bagimsizligi 
    acisindan avantaj saglamaktadir.
    
    \item \textbf{Moduler Tasarimin Degeri:} Route'larin ve 
    veritabani islemlerinin moduler yapida organize edilmesi, 
    hata ayiklama ve gelistirme sureclerini kolaylastirmistir.
    
    \item \textbf{API First Yaklasimi:} Backend API'nin once 
    tasarlanmasi ve test edilmesi, frontend gelistirmeyi 
    hizlandirmistir.
\end{enumerate}

\section{Gelecek Calismalar ve Oneriler}
\label{sec:gelecek-calismalar}

\subsection{Kisa Vadeli Gelistirmeler}
\label{subsec:kisa-vadeli}

\begin{enumerate}
    \item \textbf{Bildirim Sistemi:} Push notification ile gorev 
    hatirlatmalari ve durum guncellemelerinin personele iletilmesi.
    
    \item \textbf{Fotograf Ekleme:} Gorev tamamlama sirasinda 
    temizlik oncesi/sonrasi fotograf yukleme ozelligi.
    
    \item \textbf{QR Kod Entegrasyonu:} Lokasyonlara yerlestirilecek 
    QR kodlar ile hizli gorev isaretleme.
    
    \item \textbf{Gelismis Raporlama:} Grafik ve chart tabanli 
    gorsel raporlar ve Excel/PDF disa aktarim.
\end{enumerate}

\subsection{Orta Vadeli Gelistirmeler}
\label{subsec:orta-vadeli}

\begin{enumerate}
    \item \textbf{Mobil Uygulama:} React Native veya Flutter ile 
    yerel mobil uygulama gelistirmesi.
    
    \item \textbf{IoT Entegrasyonu:} Hareket sensorleri ve doluluk 
    sensorleri ile akilli temizlik planlamasi.
    
    \item \textbf{Yapay Zeka:} Makine ogrenmesi ile optimal 
    personel cizelgeleme ve talep tahmini.
\end{enumerate}

\section{Sonuc}
\label{sec:sonuc-bolumu}

Bu bitirme projesi kapsaminda, Kirikkale Universitesinin temizlik 
hizmetlerinin dijital ortamda yonetilebilmesi icin kapsamli bir 
web uygulamasi gelistirilmistir. Gelistirilen Kurumsal Temizlik 
Takip ve Yonetim Sistemi (KTTYS), modern web teknolojilerini 
kullanarak manuel sureclerin dijitallestirilmesini, gorevlerin 
izlenebilir hale getirilmesini ve performansin olculebilmesini 
saglamaktadir.

Proje, yazilim muhendisligi prensiplerini uygulayarak, moduler 
ve surdurulebilir bir kod tabani olusturmustur. RESTful API 
tasarimi, rol tabanli erisim kontrolu ve responsive kullanici 
arayuzu ile profesyonel bir kurumsal uygulama standartlarini 
karsilamaktadir.

Sonuc olarak, bu bitirme projesi hem teknik beceri gelistirme 
hem de gercek dunya problemlerine cozum uretme acisindan 
degerli bir ogrenme deneyimi olmustur.
